% from ega4-8-fr.tex
\oldpage[IV-3]{30}
pour $1\leq i\leq n$~; on en déduit pour tout $\nu\geq\mu$ un
homomorphisme de $E$-algèbres
$\theta_\nu=\varphi_{\nu\mu}\circ\theta_\mu$ de $G$ dans $F_\nu$, et il est
clair que $\theta$ est la limite inductive de ce système inductif
d'homomorphismes, ce qui achève de démontrer le lemme.

Passons maintenant au cas où $X_\alpha$ est quasi-compact et $Y_\alpha$
localement de type fini sur $S_\alpha$. Posons
$Z_\alpha=X_\alpha\times_{S_\alpha}Y_\alpha$ et introduisons le système
projectif correspondant des
$Z_\lambda=Z_\alpha\times_{S_\alpha}S_\lambda=
X_\lambda\times_{S_\lambda}Y_\lambda$ et sa limite
$Z=Z_\alpha\times_{S_\alpha}S=X\times_S Y$~; les bijections canoniques
(I, 3.3.14) donnent des diagrammes commutatifs
\[
\begin{tikzcd}[column sep=large,row sep=large]
\operatorname{Hom}_{S_\lambda}(X_\lambda,Y_\lambda)
  \ar[r,"e_{\mu\lambda}"] \ar[d,"\wr"] &
\operatorname{Hom}_{S_\mu}(X_\mu,Y_\mu)
  \ar[r,"e_\mu"] \ar[d,"\wr"] &
\operatorname{Hom}_S(X,Y) \ar[d,"\wr"] \\
\operatorname{Hom}_{X_\lambda}(X_\lambda,Z_\lambda)
  \ar[r,"e_{\mu\lambda}"'] &
\operatorname{Hom}_{X_\mu}(X_\mu,Z_\mu)
  \ar[r,"e_\mu"'] &
\operatorname{Hom}_X(X,Z)
\end{tikzcd}
\]
et par suite on est ramené à prouver que (8.8.1.1) est injective dans le cas
particulier où $S_\alpha=X_\alpha$ (compte tenu de (I, 3.4)). De plus, comme
$X_\alpha$ est quasi-compact, donc réunion finie d'ouverts affines, on peut
supposer $X_\alpha$ affine ($L$ étant filtrant). Supposons donc donnés deux
$X_\alpha$-morphismes $f'_\alpha:X_\alpha\to Y_\alpha$,
$f''_\alpha:X_\alpha\to Y_\alpha$ tels que $f'$ et $f''$ soient des
$X$-morphismes égaux de $X$ dans $Y$~; il s'agit de prouver que
$f'_\mu=f''_\mu$ pour $\mu\geq\alpha$ assez grand. Comme $X_\alpha$ est
quasi-compact, il en est de même de
$f'_\alpha(X_\alpha)\cup f''_\alpha(X_\alpha)$, et comme $Y_\alpha$ est
localement de type fini sur $X_\alpha$,
$f'_\alpha(X_\alpha)\cup f''_\alpha(X_\alpha)$ est contenu dans une réunion
finie d'ouverts affines $V_{i\alpha}$ de $Y_\alpha$, de type fini sur
$X_\alpha$. Posons
$U'_{i\alpha}=f_\alpha^{'-1}(V_{i\alpha})$,
$U''_{i\alpha}=f_\alpha^{''-1}(V_{i\alpha})$,
$U_{i\alpha}=U'_{i\alpha}\cap U''_{i\alpha}$,
$U_\alpha=\bigcup_iU_{i\alpha}$. L'hypothèse $f'=f''$ entraîne
$v_\alpha^{-1}(U'_{i\alpha})=v_\alpha^{-1}(U''_{i\alpha})$, ces deux ensembles
étant respectivement égaux à $f^{'-1}(w_\alpha^{-1}(V_{i\alpha}))$ et à
$f^{''-1}(w_\alpha^{-1}(V_{i\alpha}))$. Comme les $V_{i\alpha}$ recouvrent
$Y_\alpha$, on a $v_\alpha^{-1}(U_\alpha)=f^{'-1}(Y)=X$. Comme $X_\alpha$
est quasi-compact et que toute partie ouverte de $X_\lambda$ est
ind-constructible, il résulte de (8.3.4) qu'il y a un indice
$\lambda\geq\alpha$ tel que les
$U_{i\mu}=v_{\alpha\mu}^{-1}(U_{i\alpha})$ forment un \emph{recouvrement} de
$X_\mu$. Remplaçant $\alpha$ par $\mu$, on peut donc supposer que les
$U_{i\alpha}$ recouvrent $X_\alpha$~; cela entraîne que pour tout
$x\in X_\alpha$, il y a un voisinage ouvert affine $W(x)$ contenu dans un des
$U_{i\alpha}$, autrement dit tel que les restrictions de $f'_\alpha$ et
$f''_\alpha$ à $W(x)$ appliquent $W(x)$ dans un \emph{même} ouvert affine
$V_{i\alpha}$. Comme $X_\alpha$ est quasi-compact, il est recouvert par un
nombre fini d'ouverts affines $W(x_k)$~; en vertu du lemme (8.8.2.1) et du fait
que $L$ est filtrant, il existe un $\lambda\geq\alpha$ tel que les
restrictions de $f'_\lambda$ et $f''_\lambda$ à chacun des ouverts
$v_{\alpha\lambda}^{-1}(W(x_k))$ soient égales, d'où
$f'_\lambda=f''_\lambda$.

Supposons maintenant $X_\alpha$ quasi-compact et quasi-séparé et $Y_\alpha$
localement de présentation finie sur $S_\alpha$, et prouvons que (8.8.1.1) est
surjective. Supposons donc donné un $X$-morphisme $f:X\to Y$. Comme $X$ est
quasi-compact, il en est de même
